% Tutorial set: 03
% Source: T3-MDP and RL.pdf, p. 1 (PDF p. 1), Questions 3.1--3.2(a--d)
% Reference: T3-MDP and RL- solution.pdf, PDF pp. 1--2;
% T3-MDP and RL_2(A) details.pdf, pp. 1--18, especially pp. 17--18
% Session date: 2026-09-18 (Week 6)
% Verification: derivations and tabular updates checked against the stated model

\exercisegroup{SC3000-TUT03}{Tutorial 03：MDP 与 Reinforcement Learning}
\label{tutorial:SC3000-TUT03}

\exercisemeta
  {SC3000-TUT03}
  {2026-09-18（Week 6）}
  {T3-MDP and RL，p.~1，Questions 3.1--3.2；参考答案 PDF pp.~1--2；3.2(a) 补充详解 pp.~1--18}
  {建模、两轮 Value Iteration、MC 与 Q-learning 数值已核对}
  {已确认（2026-09-18）}

\exercisequestion{SC3000-TUT03-Q01}{3.1：Inventory Management 的 MDP Formulation}

\textbf{对应讲义：}\relatedtopic{SC3000-L03-T03}{MDP、Markov Property 与 Policy}。

\textbf{题目。}仓库容量为 $W$ 件。每月月初观察当前 inventory（库存），根据已知的随机 demand（需求）分布决定向供应商订购多少件。自行定义固定的单位采购成本、仓储成本与售价，以最大化利润为目标，给出 state、action、transition 与 reward 的 MDP formulation（马尔可夫决策过程建模）。

\begin{checkedsolution}
\textbf{建模约定。}每期先观察库存，再订货到货，最后满足需求；未满足需求直接流失，不产生 backorders（积压订单）。采用各月需求独立同分布、分布 $p_d=P(D_t=d)$ 已知的简化模型。以上是使模型明确的假设；若存在到货延迟、积压订单或历史相关需求，需要补充相应状态。

\textbf{State 与 action。}令 $s_t$ 为月初库存、$a_t$ 为本次订购数量，则
\[
  \boxed{S=\{0,1,\ldots,W\}},\qquad
  \boxed{A(s_t)=\{0,1,\ldots,W-s_t\}}.
\]
动作必须满足 $s_t+a_t\le W$；只列全局动作范围 $\{0,\ldots,W\}$，还不足以表达各状态下的容量限制。

\textbf{Transition。}补货后库存为 $q=s+a$。给定实际需求 $d$，销量是 $\min(d,q)$，下一期库存为
\[
  \boxed{s_{t+1}=\max\{s_t+a_t-d_t,0\}}.
\]
因此，对 $j\in S$，
\[
  \boxed{T(j\mid s,a)=
  \begin{cases}
    P(D\ge q), & j=0,\\
    P(D=q-j), & 1\le j\le q,\\
    0, & j>q.
  \end{cases}}
\]
正库存对应确定的需求量；库存为零则合并了需求恰好耗尽库存和需求超过库存的所有情况。概率满足
\[
  P(D\ge q)+\sum_{j=1}^{q}P(D=q-j)=1.
\]
$q=0$ 时求和为空，下一状态必为零。

\textbf{Reward。}设单位采购成本为 $b$、单位仓储费为 $h$、单位售价为 $v$；沿用参考答案对补货后的 $q$ 件收取仓储费，而非对期末库存收费。实际一期利润与期望奖励分别为
\[
  \boxed{r(s,a,d)=v\min(d,s+a)-ba-h(s+a)},
\]
\[
\begin{aligned}
  R(s,a)
    &=v\,\mathbb E[\min(D,q)]-ba-hq\\
    &=v\left[\sum_{d=0}^{q}d\,p_d+qP(D>q)\right]-ba-hq.
\end{aligned}
\]
\textbf{勘误：}参考答案的销售收入求和遗漏了需求超过库存时的收入；本模型需保留 $vqP(D>q)$。采购成本只针对新买的 $a$ 件，销量不超过可用库存 $q$。

原题未指定经营期限或折扣率。若采用无限期折扣目标，可另行选定 $0\le\gamma<1$，最大化 $\mathbb E_\pi[\sum_{t\ge0}\gamma^t r_t]$；这是目标函数的建模选择，并非题目给定数值。
\end{checkedsolution}

\exercisequestion{SC3000-TUT03-Q02}{3.2：Gridworld 的 Planning 与 Sample-Based Learning}

\textbf{对应讲义：}\relatedtopic{SC3000-L04-T01}{Value Iteration}；\par
\relatedtopic{SC3000-L04-T03}{Model-Free RL}；
\relatedtopic{SC3000-L04-T04}{First-Visit MC}；\par
\relatedtopic{SC3000-L05-T01}{Q-Learning}；
\relatedtopic{SC3000-L05-T02}{Off-Policy Control}。

\textbf{共同题设。}Gridworld（网格世界）有两行三列，坐标按 $(\text{row},\text{column})$ 记。Agent 从 $(1,1)$ 出发；进入 $(2,3)$ 得 $+5$，进入 $(1,3)$ 得 $-5$，随后终止；进入其他状态奖励为零。

\begin{center}
  % Original redraw from T3-MDP and RL, p. 1, Figure (a).
% Labels +5/-5 denote rewards on entry, not terminal state values.
\begin{tikzpicture}[x=2.2cm,y=1.0cm,font=\small]
  \fill[green!12] (2,1) rectangle (3,2);
  \fill[red!10] (2,0) rectangle (3,1);
  \draw[step=1,draw=noteBlue] (0,0) grid (3,2);
  \node[align=center] at (0.5,0.5) {$(1,1)$\\起点 $S$};
  \node at (1.5,0.5) {$(1,2)$};
  \node[align=center] at (2.5,0.5) {$(1,3)$\\$-5$；终止};
  \node at (0.5,1.5) {$(2,1)$};
  \node at (1.5,1.5) {$(2,2)$};
  \node[align=center] at (2.5,1.5) {$(2,3)$\\$+5$；终止};
  \draw[-{Stealth},thick] (-0.45,0.3)--(-0.45,1.6)
    node[above] {North};
  \node[below,align=center,font=\footnotesize] at (1.5,-0.12)
    {第 1 行在下，第 2 行在上；终点奖励在进入时取得};
\end{tikzpicture}

\end{center}

每步选择 North、South、West 或 East（分别简称 $N,S,W,E$）：以 $0.8$ 的概率沿所选方向移动，以各 $0.1$ 的概率沿两个垂直方向移动；撞墙则停留原格。上图按原题 Figure (a) 重绘；图内 $\pm5$ 表示进入奖励，不是 terminal value（终点后续价值）。

\begin{checkedsolution}
\nopagebreak[4]
\textbf{(a) 两轮 Value Iteration。}已知转移概率，取 $\gamma=0.9$、$V_0(s)=0$；两个终点的价值始终为零。要求计算所有状态的 $V_1,V_2$。采用 synchronous update（同步更新）：
\[
  Q_{k+1}(s,a)=\sum_{s'}T(s'\mid s,a)
  [R(s,a,s')+0.9V_k(s')],\qquad
  V_{k+1}(s)=\max_a Q_{k+1}(s,a).
\]
一轮内只能使用旧表 $V_k$，不能混入刚得到的 $V_{k+1}$。这里的 $Q_{k+1}$ 是已知模型下的动作评分，不是 (d) 中的采样 Q-learning 更新。

第一轮只有 $(2,2)$ 的最大期望即时奖励为正：
\[
  V_1(2,2)=Q_1((2,2),E)=0.8\times5=4.
\]
$(1,2)$ 选择 $W$ 可避免进入负终点，故 $V_1(1,2)=0$；其余状态第一轮也为零。两轮的完整状态值为
\begin{center}
\begin{tabular}{@{}lrrrrrr@{}}
  \toprule
  $s$ & $(1,1)$ & $(1,2)$ & $(1,3)$ & $(2,1)$ & $(2,2)$ & $(2,3)$ \\
  \midrule
  $V_0(s)$ & 0 & 0 & 0 & 0 & 0 & 0 \\
  $V_1(s)$ & 0 & 0 & 0 & 0 & 4 & 0 \\
  $V_2(s)$ & 0 & 2.38 & 0 & 2.88 & 4.36 & 0 \\
  \bottomrule
\end{tabular}
\end{center}

第二轮各非终点的四个动作评分如下，每行取最大值：
\begin{center}
\begin{tabular}{@{}lrrrrr@{}}
  \toprule
  状态 & $N$ & $S$ & $W$ & $E$ & $V_2$ \\
  \midrule
  $(1,1)$ & 0 & 0 & 0 & 0 & 0 \\
  $(1,2)$ & 2.38 & $-0.50$ & 0.36 & $-3.64$ & 2.38 \\
  $(2,1)$ & 0.36 & 0.36 & 0 & 2.88 & 2.88 \\
  $(2,2)$ & 3.38 & 0.50 & 0.36 & 4.36 & 4.36 \\
  \bottomrule
\end{tabular}
\end{center}
关键计算为
\[
\begin{aligned}
  Q_2((1,2),N)&=0.8(0.9\times4)+0.1(-5)+0.1\times0=2.38,\\
  Q_2((2,1),E)&=0.8(0.9\times4)=2.88,\\
  Q_2((2,2),E)&=0.8\times5+0.1(0.9\times4)+0.1\times0=4.36.
\end{aligned}
\]
最后一式中，向北偏移会撞墙并留在 $(2,2)$，贡献 $0.1\times0.9\times4=0.36$。\textbf{简版答案的 $4.0$ 漏算此项，补充详解 p.~18 的 $4.36$ 正确。}进入终点只计即时奖励，其后续价值为零，不能把 $\pm5$ 再作为终点价值重复计入。

同步更新使奖励信息逐轮向外传播。第二轮起点 $(1,1)$ 的一步可达状态在 $V_1$ 中仍全为零，所以 $V_2(1,1)=0$；不能提前使用邻居新算出的 $2.38$ 或 $2.88$。

\paragraph{(b) 未知转移概率时的学习条件.}
问题是：若不知道 $T(s'\mid s,a)$，为了学习 optimal policy（最优策略），agent 需要具备什么能力？它必须能够与环境交互，尝试动作，观察 reward 与 next state，获得 $(s,a,r,s')$ 经验；探索还应充分覆盖相关 state--action pairs（状态--动作对）。可以先估计模型，也可以如 Q-learning 一样直接学习动作价值，不要求显式恢复转移概率。仅拥有少数固定动作的经验，不能据此保证已学到最优策略。

\paragraph{(c) Monte Carlo 状态价值估计.}
始终选择 $E$ 的策略产生以下三条轨迹；实际向北移动是随机偏移，不代表选择了 $N$：
\[
\begin{aligned}
  \tau_1 &: (1,1)\to(1,2)\to(1,3), && r_{\mathrm{final}}=-5,\\
  \tau_2 &: (1,1)\to(1,2)\to(2,2)\to(2,3), && r_{\mathrm{final}}=+5,\\
  \tau_3 &: (1,1)\to(2,1)\to(2,2)\to(2,3), && r_{\mathrm{final}}=+5.
\end{aligned}
\]
本小题取 $\gamma=1$，要求估计 $(1,1)$ 和 $(2,2)$ 的状态价值。MC 对访问该状态后得到的完整 return（回报）取平均；未访问该状态的轨迹不计入分母，也不当成一次零回报。因此
\[
  \boxed{\widehat V(1,1)=\frac{-5+5+5}{3}=\frac53},\qquad
  \boxed{\widehat V(2,2)=\frac{5+5}{2}=5}.
\]
估计对象是上述固定策略的 $V^\pi$，不是最优价值 $V^*$；有限样本均值也不等于精确价值。本题每条轨迹没有重复访问同一状态，故 first-visit 与 every-visit MC 结果相同。

\paragraph{(d) 两次 Trial 的 Q-learning 更新.}
要求按 $\tau_1,\tau_2$ 的时间顺序逐步更新，取 $\alpha=0.1$，初始所有 $Q(s,a)=0$，trial 之间不重置。原题 (d) 未重新明确 $\gamma$；以下沿用参考答案的 $0.9$：
\[
  Q(s,a)\leftarrow Q(s,a)+0.1
  \left[r+0.9\max_{a'}Q(s',a')-Q(s,a)\right].
\]
到达终点时后续项为零；最大值比较下一状态的全部动作。更新对象是实际选择的 $E$，而不是随机实现的位移方向。

Trial 1 的两步为
\[
\begin{aligned}
  Q((1,1),E)&\leftarrow0+0.1(0+0.9\times0-0)=0,\\
  Q((1,2),E)&\leftarrow0+0.1(-5+0-0)=-0.5.
\end{aligned}
\]
Trial 2 保留上述结果，再依次更新：
\[
\begin{aligned}
  Q((1,1),E)&\leftarrow0+0.1(0+0.9\times0-0)=0,\\
  Q((1,2),E)&\leftarrow-0.5+0.1[0+0.9\times0-(-0.5)]=-0.45,\\
  Q((2,2),E)&\leftarrow0+0.1(5+0-0)=0.5.
\end{aligned}
\]
第一步虽有 $Q((1,2),E)=-0.5$，其余动作仍为零，故最大值为零；第二步发生时，$(2,2)$ 的 Q 值尚未更新。最后得到的 $0.5$ 不会自动倒传重算前两步，需要后续更新才能继续传播。其余 Q 值均保持零。

本题这五次更新中的后续最大 Q 值恰好全为零，因此即使把 (c) 的 $\gamma=1$ 延续到 (d)，这两次 trial 的数值也不变；一般情况下不能混用折扣率。
\end{checkedsolution}
